% Part: many-valued-logic
% Chapter: sequent-calculus
% Section: introduction

\documentclass[../../../include/open-logic-section]{subfiles}

\begin{document}

\olfileid{mvl}{seq}{int}

\olsection{ভূমিকা}

ধ্রুপদি যুক্তিবিদ্যার সিকোয়েন্ট কলন একটি দক্ষ ও সরল
!!{derivation} পদ্ধতি। কোনো বহুমানী যুক্তিবিদ্যা যদি সসীমসংখ্যক
সত্যমানযুক্ত ম্যাট্রিক্স দিয়ে সংজ্ঞায়িত হয়, অর্থাৎ $V$ সসীম হয়,
তবে তার জন্যও সিকোয়েন্ট কলন দেওয়া সম্ভব। কীভাবে তা করা যায়, সেই
ধারণাটি আসে ধ্রুপদি ক্ষেত্রে সিকোয়েন্টের অর্থ ও অনুমান-বিধির রূপ
বিবেচনা থেকে।

মনে করো, একটি সিকোয়েন্ট
\begin{align*}
    !A_1, \dots, !A_m & \Sequent !B_1, \dots, !B_n
\intertext{নিচের !!{formula} হিসেবে ব্যাখ্যা করা যায়:}
    (!A_1 \land \cdots \land !A_m) & \lif (!B_1 \lor \cdots \lor
!B_n)
\end{align*}
% BN-SRC-329 TARGET-CORRECTION-BEGIN
\par\smallskip\noindent\textnormal{\small\textbf{উৎস-সংশোধন (BN-SRC-329):} হিমায়িত উৎসে বাঁপাশের শেষ সূত্রটি $!A_n$ লেখা, অথচ সঙ্গে দেওয়া ব্যাখ্যায় সেটি $!A_m$। দুই পাশের দৈর্ঘ্য আলাদা হতে পারে; তাই বাঁপাশের সূচক $m$ করা হয়েছে।} % BN-SRC-329 TARGET-CORRECTION-END
অর্থাৎ, !!^a{valuation}~$\pAssign{v}$ সিকোয়েন্ট
$\Gamma \Sequent \Delta$-কে \emph{পরিতৃপ্ত করে} তখনই, যখন
কোনো $!A \in \Gamma$-র জন্য $\pValue{v}(!A) = \False$ হয়,
অথবা কোনো $!A \in \Delta$-র জন্য $\pValue{v}(!A) = \True$ হয়।
এই ব্যাখ্যায় প্রারম্ভিক সিকোয়েন্ট $!A \Sequent !A$ সর্বদা
পরিতৃপ্ত হয়, কারণ হয় $\pValue v(!A) = \True$,
নয়তো $\pValue{v}(!A) = \False$।
% BN-SRC-330 TARGET-CORRECTION-BEGIN
\par\smallskip\noindent\textnormal{\small\textbf{উৎস-সংশোধন (BN-SRC-330):} হিমায়িত উৎসে শেষ সমীকরণে $\pValue(!A)$ লিখে সত্যমান-অর্পণের বাধ্যতামূলক আর্গুমেন্ট $v$ বাদ গেছে। এখানে $\pValue{v}(!A)$ লেখা হয়েছে; একই অনুচ্ছেদের আগের সমীকরণেও এই অর্পণই ব্যবহৃত।} % BN-SRC-330 TARGET-CORRECTION-END

পাশের সূত্রগুলি $\Gamma$, $\Delta$ বাদ দিয়ে $\Log{LK}$-এ
শর্তবচনের অনুমান-বিধিগুলি নিচে দেওয়া হলো:

\begin{defish}
    \Axiom$ \fCenter !A$
    \Axiom$ !B \fCenter $
    \RightLabel{\LeftR{\lif}}
    \BinaryInf$ !A \lif !B \fCenter $
    \DisplayProof
    \hfill
    \Axiom$ !A \fCenter !B$
    \RightLabel{\RightR{\lif}}
    \UnaryInf$ \fCenter !A \lif !B $
    \DisplayProof
\end{defish}

সিকোয়েন্টের উপরোক্ত অর্থগত ব্যাখ্যা প্রয়োগ করলে
$\LeftR{\lif}$ বিধিটি বলে: $\pValue v(!A) = \True$ এবং
$\pValue v(!B) = \False$ হলে $\pValue v(!A \lif !B) = \False$।
অনুরূপভাবে, $\RightR{\lif}$ বিধিটি বলে: $\pValue v(!A) = \False$
অথবা $\pValue v(!B) = \True$ হলে $\pValue v(!A \lif !B) =
\True$। বস্তুত, এই শর্তগুলি উভমুখী শর্ত। $\LeftR{\land}$ ও
$\RightR{\lor}$ বিধির প্রচলিত রূপে সংশ্লিষ্ট শর্তগুলি উভমুখী
হতো না। তবে এই বিধিগুলির বিকল্প রূপ আছে, যাতে সেগুলি উভমুখী হয়:

\begin{defish}
    \Axiom$!A, !B, \Gamma \fCenter \Delta$
    \RightLabel{\LeftR{\land}}
    \UnaryInf$!A \land !B, \Gamma \fCenter \Delta$
    \DisplayProof
    \hfill
    \Axiom$ \Gamma \fCenter \Delta, !A, !B$
    \RightLabel{\RightR{\lor}}
    \UnaryInf$ \Gamma \fCenter \Delta, !A \lor !B$
    \DisplayProof
\end{defish}

এই মূল ধারণাটি $n$-মানী যুক্তিবিদ্যায় প্রয়োগ করলে দুই পাশের বদলে
$n$-টি অবস্থানযুক্ত সিকোয়েন্ট কলন পাওয়া যায়, প্রতি সত্যমানের জন্য
একটি অবস্থান। $V = \{\False, \Undef, \True\}$-যুক্ত ত্রিমানী
যুক্তিবিদ্যায় একটি সিকোয়েন্টের রূপ $\Gamma \mid \Pi \mid \Delta$।
এটি !!a{valuation}~$\pAssign v$-তে পরিতৃপ্ত হয় তখনই, যখন
কোনো $!A \in \Gamma$-র মান $\pValue{v}(!A) = \False$, অথবা
কোনো $!A \in \Delta$-র মান $\pValue{v}(!A) = \True$, অথবা
কোনো $!A \in \Pi$-র মান $\pValue{v}(!A) = \Undef$।
ফলে প্রারম্ভিক সিকোয়েন্ট $!A \mid !A \mid !A$ সর্বদা পরিতৃপ্ত হয়।

\end{document}
